\documentclass[a4paper,11pt]{article}
\usepackage{exam-style-341}
\examinehead{341 农业综合知识三 · 网络技术分章练习}{第07章 运输层}
\begin{document}

\unitcover{341 农业综合知识三 · 网络技术分章练习}{第07章\quad 运输层}{练习（共 100 分）}{约 75 分钟}{运输层功能与端口、套接字；UDP 的特点与首部；TCP 的特点与首部；连接管理（三次握手、四次挥手、TIME_WAIT）；序号与确认号推算；滑动窗口与流量控制；拥塞控制（慢开始、拥塞避免、快重传、快恢复）。}

\begin{knowmap}
\textbf{运输层：}为应用进程提供端到端的逻辑通信；用端口号标识进程，IP 地址＋端口号＝套接字\\[2pt]
\textbf{端口范围：}0—65535；熟知端口 0—1023（如 FTP 21、SMTP 25、HTTP 80）· 登记端口 1024—49151 · 短暂端口 49152—65535\\[2pt]
\textbf{UDP：}无连接、尽最大努力交付、面向报文、无拥塞控制；首部仅 8 字节（源端口、目的端口、长度、检验和）\\[2pt]
\textbf{TCP：}面向连接、可靠、字节流、全双工；首部固定 20 字节，含序号、确认号、窗口等字段\\[2pt]
\textbf{连接管理：}三次握手（SYN $\rightarrow$ SYN+ACK $\rightarrow$ ACK）；四次挥手（FIN $\rightarrow$ ACK $\rightarrow$ FIN $\rightarrow$ ACK）；TIME\_WAIT 等待 2MSL\\[2pt]
\textbf{可靠传输与窗口：}序号＝本段第一个字节编号；确认号＝期望收到的下一个字节编号；发送窗口 $=\min(rwnd,cwnd)$；慢开始加倍、拥塞避免加 1、快重传不等超时、快恢复减半后直接拥塞避免
\end{knowmap}

\begin{keyterms}
\textbf{端口}：运输层用 16 位端口号标识主机中的应用进程，范围 0—65535。\\
\textbf{套接字}：IP 地址与端口号的组合，唯一标识互联网中的一个通信端点。\\
\textbf{UDP}：用户数据报协议，无连接、尽最大努力交付、面向报文，首部 8 字节。\\
\textbf{TCP}：传输控制协议，面向连接、可靠、字节流、全双工，首部固定 20 字节。\\
\textbf{序号}：TCP 报文段首部字段，指出本报文段数据第一个字节的编号；SYN、FIN 各占一个序号。\\
\textbf{确认号}：期望收到的下一个字节的序号，表示该序号之前的数据已全部正确收到。\\
\textbf{三次握手}：TCP 建立连接时交换 SYN、SYN+ACK、ACK 三个报文段，同步双方序号并确认收发能力。\\
\textbf{四次挥手}：TCP 释放连接时双方各发送 FIN 并各确认一次，共四个报文段。\\
\textbf{TIME\_WAIT}：主动关闭方发送最后一个 ACK 后进入的状态，等待 2MSL 后才关闭连接。\\
\textbf{滑动窗口}：TCP 以字节为单位的发送与接收窗口机制，实现可靠传输与流量控制。\\
\textbf{流量控制}：接收方通过通告接收窗口 rwnd 限制发送方发送速率，防止接收方缓冲区溢出。\\
\textbf{拥塞控制}：发送方根据网络拥塞程度调整拥塞窗口 cwnd，防止过多数据注入网络。\\
\textbf{慢开始与拥塞避免}：慢开始阶段 cwnd 每 RTT 加倍，达到 ssthresh 后转入拥塞避免，每 RTT 加 1。\\
\textbf{快重传与快恢复}：收到 3 个重复确认即立即重传丢失报文段；快恢复将 ssthresh 减半后直接进入拥塞避免。
\end{keyterms}

\begin{contrast}{UDP / TCP（每行依次为：连接与可靠性·首部与开销·典型应用）}
\citem{UDP}{无连接、尽最大努力交付、不保证可靠}{首部 8 字节，开销小}{DNS、DHCP、实时音频视频、广播与多播}
\citem{TCP}{面向连接、可靠、字节流、全双工}{首部固定 20 字节，含序号、确认号、窗口}{HTTP、FTP、SMTP、POP3、远程登录}
\end{contrast}

\begin{contrast}{流量控制 / 拥塞控制（每行依次为：控制目的·控制变量·主要机制）}
\citem{流量控制}{防止发送方发得太快导致接收方缓冲区溢出}{接收方通告的接收窗口 rwnd}{接收方通过窗口字段动态通告可用缓冲区大小}
\citem{拥塞控制}{防止过多数据注入网络导致网络性能下降}{发送方维护的拥塞窗口 cwnd}{慢开始、拥塞避免、快重传、快恢复}
\end{contrast}

\begin{contrast}{三次握手 / 四次挥手（每行依次为：目的·报文段序列·为什么是这个次数）}
\citem{三次握手}{建立连接、同步双方初始序号、确认收发能力}{SYN $\rightarrow$ SYN+ACK $\rightarrow$ ACK}{两次握手无法确认客户端收到了服务器的 SYN+ACK，且可能因失效请求建立错误连接}
\citem{四次挥手}{释放连接，双方各自关闭发送方向}{FIN $\rightarrow$ ACK $\rightarrow$ FIN $\rightarrow$ ACK}{TCP 是全双工的，一方发送 FIN 只表示自己不再发送数据，另一方可能还有数据要发，故 FIN 与 ACK 不能合并}
\end{contrast}

\qtypename{一、单项选择题}{10 题，每题 3 分，共 30 分}
\q 运输层用（\quad）来标识主机中的应用进程。
\choicesline{A. IP 地址}{B. MAC 地址}{C. 端口号}{D. 域名}

\q UDP 用户数据报的首部长度是（\quad）
\choicesline{A. 4 字节}{B. 8 字节}{C. 20 字节}{D. 40 字节}

\q TCP 报文段首部的固定部分长度是（\quad）
\choicesline{A. 8 字节}{B. 20 字节}{C. 40 字节}{D. 60 字节}

\q 下列端口号中，属于熟知端口（系统端口）范围的是（\quad）
\choicesline{A. 80}{B. 1024}{C. 50000}{D. 65535}

\q TCP 三次握手中，服务器发送的第二个报文段的标志位是（\quad）
\choicesline{A. SYN=1，ACK=1}{B. SYN=1，ACK=0}{C. SYN=0，ACK=1}{D. FIN=1，ACK=1}

\q TCP 连接释放时，主动关闭方发送最后一个 ACK 后进入（\quad）状态，等待 2MSL。
\choicesline{A. CLOSE\_WAIT}{B. LAST\_ACK}{C. TIME\_WAIT}{D. FIN\_WAIT\_1}

\q TCP 发送方的发送窗口大小等于（\quad）
\choices{A. $\max(rwnd,cwnd)$}{B. $\min(rwnd,cwnd)$}{C. $rwnd+cwnd$}{D. $rwnd$}

\q 在 TCP 慢开始阶段，若未发生拥塞，拥塞窗口 cwnd 每个 RTT（\quad）
\choicesline{A. 加 1}{B. 加倍}{C. 减半}{D. 不变}

\q TCP 快重传算法的触发条件是（\quad）
\choices{A. 收到 3 个重复确认}{B. 超时计时器到期}{C. 收到 1 个重复确认}{D. 接收窗口变为 0}

\q TCP 报文段首部中的确认号字段表示（\quad）
\choices{A. 本报文段已发送的字节数}{B. 期望收到的下一个字节的序号}
{C. 本报文段第一个字节的序号}{D. 接收窗口的字节数}

\qtypename{二、多项选择题}{3 题，每题 4 分，共 12 分；每题有两个或两个以上正确项}
\q TCP 的主要特点包括（\quad）
\choices{A. 面向连接}{B. 可靠交付}{C. 面向字节流}{D. 无连接}

\q UDP 的主要特点包括（\quad）
\choices{A. 无连接}{B. 首部开销小}{C. 尽最大努力交付}{D. 提供可靠交付}

\q TCP 拥塞控制使用的方法有（\quad）
\choicesline{A. 慢开始}{B. 拥塞避免}{C. 快重传}{D. 流量控制}

\qtypename{三、判断题}{6 题，每题 2 分，共 12 分；正确的画“√”，错误的画“×”并写出改正后的表述}
\q TCP 首部固定部分为 20 字节，UDP 首部为 8 字节。\quad（\quad）\quad 改正：\fixblank

\q TCP 报文段的确认号为 500，表示接收方已经收到了 500 字节的数据。\quad（\quad）\quad 改正：\fixblank

\q 三次握手可以防止失效的连接请求报文段突然又传送到服务器而建立错误的连接。\quad（\quad）\quad 改正：\fixblank

\q 拥塞窗口 cwnd 越大越好，越大说明传输效率越高。\quad（\quad）\quad 改正：\fixblank

\q 流量控制是接收方通过通告接收窗口 rwnd 来限制发送方的发送速率。\quad（\quad）\quad 改正：\fixblank

\q TCP 连接释放需要四次挥手，这与 TCP 是全双工协议有关。\quad（\quad）\quad 改正：\fixblank

\qtypename{四、计算题}{1 题，共 6 分；要求写出序号与确认号的推导过程}
\q 主机 A 与主机 B 建立 TCP 连接，A 的初始序号 ISN=1000（SYN 报文段），B 的初始序号 ISN=800。请回答：
\begin{enumerate}[leftmargin=2em, itemsep=1pt, topsep=2pt]
  \item 三次握手中 A 的 SYN 报文段占用哪个序号？B 的 SYN+ACK 报文段的序号和确认号分别是多少？握手完成后 A 下一次发送数据的序号是多少？
  \item A 发送序号为 1001、携带 300 字节数据的报文段，B 回复的确认号是多少？
  \item A 接着发送序号为 1301、携带 500 字节数据的报文段，B 回复的确认号是多少？
  \item 若 A 随后发送 FIN 报文段（序号 1801），B 回复的确认号是多少？
\end{enumerate}
\anslines{5cm}

\newpage
\qtypename{五、简答题}{2 题，每题 10 分，共 20 分}
\q 比较 TCP 与 UDP 的主要区别（从连接、可靠性、首部、流量控制与拥塞控制、适用场合等方面作答）。\anslines{5cm}

\q 简述 TCP 三次握手与四次挥手的过程，并说明为什么建立连接需要三次、释放连接需要四次，以及 TIME\_WAIT 状态的作用。\anslines{5cm}

\qtypename{六、综合题}{1 题，共 20 分}
\q 某 TCP 连接的 MSS=1000 字节，初始拥塞窗口 cwnd=1 MSS，初始慢开始门限 ssthresh=16 MSS。请回答：
\begin{enumerate}[leftmargin=2em, itemsep=1pt, topsep=2pt]
  \item 列出第 1—8 个 RTT 结束时 cwnd 的值，并说明各阶段使用的拥塞控制方法（5 分）。
  \item 若在第 5 个 RTT 结束后发生超时（此时 cwnd=16 MSS），求新的 ssthresh 和 cwnd，并列出之后 5 个 RTT 结束时 cwnd 的值（5 分）。
  \item 若某时刻 cwnd=4000 字节，接收方通告的接收窗口 rwnd=3000 字节，求发送窗口；此时已发送但尚未确认的数据为 1000 字节，还能发送多少字节（4 分）。
  \item 说明快重传和快恢复的基本思想（6 分）。
\end{enumerate}
\anslines{8cm}

\end{document}
